The hydrostatic equilibrium inside the star means that in each point of the
inside of the star the gravitational forces are compensated by the hydrostatic
pressure forces. That means that the mater of the star remains localized in a
region of space.

The total pressure of the stellar gas has two components: the pressure due to
the movement of the stellar-gas particles $\left(p_{\mathrm{gaz}}\right)$ and
the pressure due to the emitted radiation by the stellar-gas particles
$\left(p_{\mathrm{rad}}\right)$, thus:
\begin{align*}
  p_{\mathrm{total}} &= p_{\mathrm{gaz}} + p_{\mathrm{rad}}; \\
  p_{\mathrm{gaz}} &= p_{\mathrm{H}} + p_{\mathrm{He}}; \\
  p_{\mathrm{rad}} &= \frac{1}{3}aT; % sic: source writes aT here and below (should be aT^4)
\end{align*}
\[
  p_{\mathrm{total}}
  = \rho\,\frac{n\mu_{\mathrm{He}} + (1-n)\mu_{\mathrm{H}}}{\mu_{\mathrm{H}}\mu_{\mathrm{He}}}\,RT
    + \frac{1}{3}aT. \tag{1}
\]
In order to calculate the gravitational pressure of the stellar gas let's
consider a narrow cylinder with section area $\Delta S$, along the radius of
the star. See the figure below. If the total gravitational force acting on this
cylinder is $\vec{F}_g$ than the gravitational pressure exert by the
gas-column is $p_{\mathrm{grav}} = F_g/\Delta S$.
% FIGURE NEEDED: cross-section of the star with the narrow radial gas cylinder of section Delta S and the forces F_1, F_2, F_g and radiation pressures marked (2014 TH solutions, pp. 19-20, Fig. 15a)

In order to calculate $\vec{F}_g$ let divide the cylinder in $n$ identically
small cylinders each of it with height $\Delta r$ and mass $\Delta m$.
Considering an homogenous star:
\begin{align*}
  F_1 &= K\frac{\Delta m \cdot M_1}{\left(r - \dfrac{\Delta r}{2}\right)^{2}}
       = K\frac{\Delta m \cdot M_1}{r^2\left(1 - \dfrac{\Delta r}{2r}\right)^{2}}
       = K\frac{\Delta m \cdot M_1}{r^2}\left(1 - \frac{\Delta r}{2r}\right)^{\!-2}; \\
  \frac{M}{r^3} &= \frac{M_1}{(r-\Delta r)^3}; \quad
  M_1 = M\cdot\frac{(r-\Delta r)^3}{r^3} = M\cdot\left(1 - \frac{\Delta r}{r}\right)^{\!3}; \\
  F_1 &\approx K\frac{\Delta m \cdot M}{r^2}\left(1 - 2\frac{\Delta r}{r}\right); \\
  F_2 &\approx K\frac{\Delta m \cdot M}{r^2}\left(1 - 3\frac{\Delta r}{r}\right); \\
  &\dotso \\
  F_n &\approx K\frac{\Delta m \cdot M}{r^2}\left(1 - (n+1)\frac{\Delta r}{r}\right); \\
  F_g &= F_1 + F_2 + \dotsb + F_n; \\
  F_g &= K\frac{m\cdot M}{n r^2}\left[n - \frac{1}{n}\,\frac{(1+n)n}{2}\right]; \quad
  F_g = K\frac{m\cdot M}{r^2}\left[1 - \frac{(1+n)}{2n}\right]; \\
  F_g &= K\frac{m\cdot M}{r^2}\,\frac{n-1}{2n}; \quad
  \frac{n-1}{2n} \approx \frac{1}{2}; \\
  F_g &= K\frac{m\cdot M}{2r^2}. \tag{2}
\end{align*}
Thus the gravitational pressure is
\begin{align*}
  p_g &= \frac{F_g}{\Delta S}
       = \frac{1}{\Delta S}\,
         \frac{\rho\cdot r\cdot\Delta S\cdot\rho\cdot\dfrac{4\pi r^3}{3}}{2r^2} \\
  p_g &= \frac{2}{3}\pi K r^2 \rho^2
\end{align*}
Using the relations (1) and (2) in the pressures equilibrium relationship
\[
  p_{\mathrm{total}} = p_{\mathrm{gravitational}}
\]
Results:
\begin{align*}
  \frac{2}{3}\pi K r^2\rho^2
    &= \rho\,\frac{n\mu_{\mathrm{He}} + (1-n)\mu_{\mathrm{H}}}{\mu_{\mathrm{H}}\mu_{\mathrm{He}}}\,RT
       + \frac{1}{3}aT; \\
  \frac{2}{3}\pi K r^2\cdot\rho^2
    - \frac{n\mu_{\mathrm{He}} + (1-n)\mu_{\mathrm{H}}}{\mu_{\mathrm{H}}\mu_{\mathrm{He}}}\,RT\cdot\rho
    - \frac{1}{3}aT &= 0;
\end{align*}
This is an second degree equation in $\rho$
\[
  A\rho^2 - B\rho - C = 0
\]
Where the coefficients are
\begin{align*}
  A &= \frac{2}{3}\pi K r^2 \\
  B &= \frac{n\mu_{\mathrm{He}} + (1-n)\mu_{\mathrm{H}}}{\mu_{\mathrm{H}}\mu_{\mathrm{He}}}\,RT \\
  C &= \frac{1}{3}aT
\end{align*}
The positive solution is the valid one i.e.
\[
  \rho = \frac{B + \sqrt{B^2 + 4AC}}{2A}, \qquad
  \rho = \frac{\dfrac{n\mu_{\mathrm{He}} + (1-n)\mu_{\mathrm{H}}}{\mu_{\mathrm{H}}\mu_{\mathrm{He}}}\,RT
    + \sqrt{\left(\dfrac{n\mu_{\mathrm{He}} + (1-n)\mu_{\mathrm{H}}}{\mu_{\mathrm{H}}\mu_{\mathrm{He}}}\,RT\right)^{\!2}
    + \dfrac{8}{9}\pi K r^2 aT}}
    {\dfrac{4}{3}\pi K r^2}
\]
